% !TeX root = ../main.tex

\newlecture[Weak Convergence, Skorokhod and Portmanteau]

\begin{definition}[Weak Convergence / Convergence in Distribution][7.1]
    A sequence of probability distributions $\mu_n$ with corresponding cumulative distribution functions $F_n$ converges weakly (in distribution) to a probability distribution $\mu$ with CDF $F$ if for all continuity points $x$ of $F$
    \[\lim_{n\to\infty} F_n(x) = F(x)\]
    and it is denoted by $X_n \overset{D}{\longrightarrow} X$ if $X_n \sim F_n$ and $X \sim F$, or $F_n \Rightarrow F$, or $\mu_n \Rightarrow \mu$.
\end{definition}
Notes:
\begin{enumerate}
    \item Pointwise convergence is enough: $F_n(x) \to F(x)$ when $n \to \infty$ for all continuity points $x$ of $F$.
    \item Using RV notation, $X_n \overset{D}{\longrightarrow} X$; using CDF notation, $F_n \Rightarrow F$; or using distribution notation, $\mu_n \Rightarrow \mu$.
\end{enumerate}

\begin{example}[][\relax]
    $X \sim \text{Bernoulli}(0.5)$: $X$ tossing a coin,
    \[X = \begin{cases} 1 & H\\ 0 & T\end{cases}\]
    $Y = 1 - X \sim \text{Bernoulli}(0.5)$: $Y$ tossing the same coin,
    \[Y = \begin{cases} 1 & T\\ 0 & H\end{cases}\]
    \[X \overset{D}{\longrightarrow} Y \quad \text{but} \quad Y = 1 - X\]
\end{example}

\begin{theorem}[Scheff\'e's Theorem][7.1.1]
    Let $X_n$ each have density $f_n$ and $X$ have density $f$. Then if $f_n(x) \to f(x)$ for all $x \in \R$,
    \[X_n \overset{D}{\longrightarrow} X\]
    Note:
    \begin{align*}
        \text{Def 7.1:} & \quad F_n \to F \Rightarrow X_n \overset{D}{\longrightarrow} X\\
        \text{Scheff\'e Thm:} & \quad f_n \to f \Rightarrow X_n \overset{D}{\longrightarrow} X
    \end{align*}
    \begin{proof}
        \[\begin{array}{ccc}
            F_n(x) = \displaystyle\int_{-\infty}^x f_n(u)\,du & & \\
            \downarrow & & \downarrow\\
            F(x) & \overset{\text{may not}}{\Leftarrow} & f(u)\,du
        \end{array}\]
        We need the condition that $f_n(u)$ must be either bounded / dominated by some other real-valued function, in order to say $F_n(x) \to F(x)$.

        Observe that the existence of a density implies $F$ is continuous everywhere.
        \begin{align*}
            |F_n(y) - F(y)| &= \left|\int_{-\infty}^y [f_n(x) - f(x)]\,dx\right|\\
            &\leq \int_{-\infty}^y |f_n(x) - f(x)|\,dx\\
            &\leq \int_\R |f_n(x) - f(x)|\,dx\\
            &= \int_\R [f_n(x) - f(x)]^+\,dx + \int_\R [f_n(x) - f(x)]^-\,dx\\
            &= 2\int_\R [f_n(x) - f(x)]^+\,dx
        \end{align*}
        (the last step because $\int_\R f_n = \int_\R f = 1$, so $\int_\R [f_n - f]^+ = \int_\R [f_n - f]^-$.)
        \\Then, observe that
        \[\underbrace{[f_n(x) - f(x)]^+}_{\text{dominated by } f(x)} \leq [f(x)]^+ = f(x)\]
        We have shown that inside the integral, $[f_n(x) - f(x)]^+$ is dominated by some function (and $\to 0$ pointwise), therefore by the Dominated Convergence Theorem,
        \[F_n(y) - F(y) \to 0 \quad \text{when } n \to \infty\]
    \end{proof}
\end{theorem}

\begin{theorem}[Skorokhod Representation][7.2]
    If $F_n \Rightarrow F$ (which means for any RV $X_n \sim F_n$, $X_n \overset{D}{\longrightarrow} X$, where $X \sim F$), there exist an $\Omega$ and corresponding $P$ along with random variables $Y_1, \dots, Y_n, \dots$ defined on $\Omega$, such that $Y_n \sim F_n$, $Y \sim F$ and $Y_n \overset{a.s.}{\longrightarrow} Y$.
\end{theorem}
Note: $F_n \Rightarrow F$ means
\\For \textbf{ALL} $X_n \sim F_n$, $X \sim F$: $X_n \overset{D}{\longrightarrow} X$.
\\{\color{red!70!black}Among all random variables satisfying $F_n \Rightarrow F$, there \underline{exist} a subset $Y_n \sim F_n$, $Y \sim F$, $Y_n \overset{a.s.}{\longrightarrow} Y$.}
\begin{proof}
    As long as I can find an example, that's enough to prove.
    \\Let $\Omega = [0,1]$ and $P$ be the uniform distribution on $[0,1]$.
    \[X_n(\omega) = F_n^{-1}(\omega) \quad \text{and} \quad X(\omega) = F^{-1}(\omega)\]
    That's an example.

    Note: among $X_n \overset{D}{\longrightarrow} X$, there \underline{exist} $Y_n \overset{a.s.}{\longrightarrow} Y$, where $Y_n \sim X_n$, $Y \sim X$. Finding an example is enough.
    \\Previously, we learned that we can create a random variable with any arbitrary distribution by transforming the uniform one: for $U \sim \text{Unif}(0,1)$, let $Y = F^{-1}(U)$, where $F$ is the CDF of any arbitrary random variable; then $Y \sim F$.
    \[\begin{array}{ccc}
        X_n & \overset{D}{\longrightarrow} & X\\
        \wr & & \wr\\
        F_n & \Rightarrow & F
    \end{array}\]
    Let $U \sim \text{Unif}(0,1)$,
    \[Y_n = F_n^{-1}(U)\]
    then $Y_n$ is a sequence of RV's each with CDF $F_n$.
    \[Y = F^{-1}(U)\]
    then $Y$ is a RV with CDF $F$.
    \[Y_n \overset{D}{\longrightarrow} Y \quad \text{and} \quad Y_n \overset{a.s.}{\longrightarrow} Y\]
    $Y_n$ and $Y$ is the example.
\end{proof}

\begin{theorem}[Portmanteau Lemma][7.3]
    $F_n \Rightarrow F$ \underline{if and only if} for all bounded and continuous $h: \R \to \R$, when $X_n \sim F_n$ and $X \sim F$, then
    \[Eh(X_n) \to Eh(X)\]
\end{theorem}
Remark: the Portmanteau Lemma says
\[X_n \overset{D}{\longrightarrow} X \quad \underline{\text{is equivalent to}} \quad Eh(X_n) \to Eh(X) \text{ for any } h\]
\begin{proof}
    ``If and only if'' $\Leftrightarrow$ both directions.
    \begin{enumerate}
        \item The forward direction:
        \[X_n \overset{D}{\longrightarrow} X \Rightarrow E[h(X_n)] \to E[h(X)] \quad \forall h\]
        Using the Skorokhod Representation, we choose $X_n \sim F_n$ and $X \sim F$ where $X_n \overset{a.s.}{\longrightarrow} X$. Then easily (by the continuous mapping theorem and BCT, as $h$ is bounded and continuous),
        \[E[h(X_n)] \to E[h(X)]\]
        Note: for $X_n \overset{D}{\longrightarrow} X$, using the Skorokhod representation we choose $Y_n \sim X_n$ and $Y \sim X$ $\Rightarrow$ $E[h(Y_n)] \to E[h(Y)]$. For the other random variables, even though they may not satisfy the almost sure convergence part, they share the same expected values, so the result can be generalized to all.
        \\Step 1. Using Skorokhod representation, I intentionally pick $Y_n \sim F_n$, $Y \sim F$ and $Y_n \overset{a.s.}{\longrightarrow} Y$, and I can show $E[h(Y_n)] \to E[h(Y)]$.
        \\Step 2. For any general random variables $X_n$ and $X$, we don't have $X_n \overset{a.s.}{\longrightarrow} X$, but
        \begin{align*}
            E[h(X_n)] &= E[h(Y_n)] & &\text{because } X_n \sim Y_n\\
            E[h(X)] &= E[h(Y)] & &\text{because } X \sim Y
        \end{align*}
        So $E[h(Y_n)] \to E[h(Y)]$ can be generalized to $E[h(X_n)] \to E[h(X)]$.

        \item The backward direction:
        \[E[h(X_n)] \to E[h(X)] \text{ for any } h(\cdot) \Rightarrow X_n \overset{D}{\longrightarrow} X \ / \ F_n \Rightarrow F\]
        for $X_n \sim F_n$, $X \sim F$.
        We know $E[\cdot] \to E[\cdot]$; we want to show the corresponding $F_n \to F$.
        \\Using indicator functions, we intentionally let each
        \[h_a(x) = I(x \leq a)\]
        \begin{align*}
            E[h(X_n)] &\to E[h(X)] \quad \text{for \underline{any} } h(\cdot)\\
            E[h_a(X_n)] &\to E[h_a(X)] \quad \text{for any } a
        \end{align*}
        \begin{align*}
            E[h_a(X_n)] &= E[I(X_n \leq a)] = P(X_n \leq a) = F_n(a)\\
            E[h_a(X)] &= F(a)
        \end{align*}
        \[E[h_a(X_n)] \to E[h_a(X)] \text{ for any } a \Rightarrow F_n(a) \to F(a) \text{ for any } a \Leftrightarrow X_n \overset{D}{\longrightarrow} X\]
        Note that this is just a partial proof, where we only consider the case that the CDFs are always absolutely continuous. For the complete proof, please refer to the 1st recommended textbook.
    \end{enumerate}
\end{proof}

\begin{corollary}[Continuous Mapping Theorem under Weak Convergence][7.3.1]
    If $X_n \overset{D}{\longrightarrow} X$, then $f(X_n) \overset{D}{\longrightarrow} f(X)$ for any continuous function $f$.
    \begin{proof}
        For any continuous and bounded function $g$, $g \circ f$ is also continuous and bounded. By Portmanteau's Lemma,
        \[X_n \overset{D}{\longrightarrow} X \Rightarrow E[h(X_n)] \to E[h(X)] \quad \forall h\]
        Let $h(\cdot) = g \circ f$:
        \[E[g(\fcolorbox{red!70!black}{white}{$f(X_n)$})] \to E[g(\fcolorbox{red!70!black}{white}{$f(X)$})] \quad \forall g\]
        \[\Downarrow\]
        \[f(X_n) \overset{D}{\longrightarrow} f(X)\]
    \end{proof}
\end{corollary}

\begin{corollary}[][7.3.2]
    If $X_n \overset{P}{\longrightarrow} X$, then $X_n \overset{D}{\longrightarrow} X$, but the reverse is NOT true.
\end{corollary}
